\section{Dimensional interpolation target}
\label{app:dimensional-interpolation}
% Status: open. Records the D=11/10/9 and colourless D=7/6/5 interpolation target.

This appendix records a dimensional version of the Higgs-sector problem.  It is
a theorem target, added here because it gives a compact way to organize the
Kaluza--Klein, string, Higgs, and top-quark questions.

\subsection{Full-gauge count}
\label{app:full-gauge-count}

The full-gauge count keeps colour inside the compactification symmetry.  Witten's
realistic Kaluza--Klein analysis gives the source-backed endpoint: a compact
space with \(SU(3)\times SU(2)\times U(1)\) symmetry requires seven extra
dimensions in the class considered there, giving total dimension \(D=11\)
\cite{WittenKK1981}.  The same source records the associated fermion
quantum-number obstruction for that seven-dimensional compactification.

The proposed interpolation is
\begin{equation}
\begin{array}{ccl}
D=11,\quad n_{\rm KK}=7
&:&
SU(3)\times SU(2)\times U(1)
\quad\hbox{massless-Higgs or unbroken endpoint},
\\[3pt]
D=10,\quad n_{\rm KK}=6
&:&
\hbox{DeVries interpolation sector},
\\[3pt]
D=9,\quad n_{\rm KK}=5
&:&
SU(3)\times U(1)_{\rm em}
\quad\hbox{infinitely broken endpoint}.
\end{array}
\label{eq:full-gauge-dimensional-chain}
\end{equation}
The source status of the three lines is uneven, and the distinction is part of
the claim.  The source-backed data and reconstruction targets are:
\begin{center}
\small
\begin{tabular}{p{0.20\linewidth}p{0.33\linewidth}p{0.35\linewidth}}
\toprule
Line & Source-backed content & DeVries obligation\\
\midrule
\(D=11,\ n_{\rm KK}=7\) & Witten's \(SU(3)\times SU(2)\times U(1)\)
KK endpoint and its fermion quantum-number obstruction & define the
massless-Higgs endpoint and its relation to the electroweak ray\\
\(D=10,\ n_{\rm KK}=6\) & reconstruction target & identify the
six-dimensional internal object, its light two-channel sector, and its kernel\\
\(D=9,\ n_{\rm KK}=5\) & Witten's five-sphere discussion contains an
\(SU(3)\times U(1)\) symmetry address & derive the electroweak embedding
that identifies the surviving factor with \(U(1)_{\rm em}\)\\
\bottomrule
\end{tabular}
\end{center}
The \(D=10\) middle line is a Schur-complement reconstruction label.  A
derivation must identify the six-dimensional internal object, its boundary or
compact operator, the light two-channel sector, the ordered electroweak
assignment, and the pole-matching rule.

\subsection{Colourless electroweak count}
\label{app:colourless-count}

There is a reduced count in which colour is treated as a spectator sector.  The
electroweak part of the endpoint symmetry then changes as
\begin{equation}
\begin{array}{ccl}
D=7,\quad n_{\rm KK}^{\rm ew}=3
&:&
SU(2)\times U(1),
\\[3pt]
D=6,\quad n_{\rm KK}^{\rm ew}=2
&:&
\hbox{DeVries electroweak interpolation sector},
\\[3pt]
D=5,\quad n_{\rm KK}^{\rm ew}=1
&:&
U(1)_{\rm em}.
\end{array}
\label{eq:colourless-dimensional-chain}
\end{equation}
This version gives the six-dimensional middle line a direct string-theory
source address.  Six-dimensional superstring vacua and dualities are a
standard part of the source landscape: type IIA compactified on K3 and the
heterotic string compactified on \(T^4\) provide a six-dimensional string-string
duality setting \cite{Strominger1995,Sen1995}; heterotic small instantons give
nonperturbative six-dimensional dynamics \cite{WittenSmallInstantons1995}; and
six-dimensional anomaly and string-universality constraints provide a modern
consistency framework \cite{KumarTaylor2009}.  These sources support the
existence and rigidity of six-dimensional string settings.  The colourless
\(D=6\) line is a Schur-complement reconstruction label until a source theory
derives the DeVries kernel, ordered \(J\)-assignment, and pole placement.

\subsection{Interpolation parameter}
\label{app:interpolation-parameter}

The interpolation parameters must be stated in field-theory and geometric
terms.  The electroweak field theory supplies a radial order parameter
\(t_{\rm EW}\) through the Higgs vacuum expectation value.  The
compactification descriptions supply a separate coordinate \(t_{\rm dim}\)
through radii, boundary conditions, brane-localized terms, or singularity
deformations.  A source derivation must supply a map
\begin{equation}
  \chi:\ t_{\rm dim}\longmapsto t_{\rm EW}
  \label{eq:dimensional-ray-map}
\end{equation}
or a joint source variable \(u\mapsto(t_{\rm dim}(u),t_{\rm EW}(u))\).  The
dimensional family has source operators
\begin{equation}
  \mathcal O_J(t_{\rm dim}),
  \qquad
  t_{\rm dim}=0:\ SU(3)\times SU(2)\times U(1),
  \qquad
  t_{\rm dim}=\infty:\ SU(3)\times U(1)_{\rm em},
  \label{eq:dimensional-interpolation-family}
\end{equation}
or the corresponding colourless electroweak endpoints.  The
\(U(1)_{\rm em}\) identification in the last entry is an electroweak embedding
obligation.  The middle-dimensional data required at a DeVries point are
\begin{equation}
  \mathfrak I_J(t_{\rm dim,\star})
  =
  \big(
  \mathcal B(t_{\rm dim,\star}),\,
  \mathcal H_J(t_{\rm dim,\star}),\,
  K_J(t_{\rm dim,\star},\lambda),\,
  \mathcal R_{\rm pole}
  \big),
  \label{eq:middle-dimensional-data}
\end{equation}
where \(\mathcal B(t_{\rm dim,\star})\) is the compact, boundary, or singular object,
\(\mathcal H_J(t_{\rm dim,\star})=\operatorname{span}\{h_J,a_J\}\) is the light
order-parameter/current subspace, \(K_J\) is the reduced kernel, and
\(\mathcal R_{\rm pole}\) is the matching rule to the transverse pole
propagator.  At the DeVries point the light two-channel reduction must take the
form
\begin{equation}
  \mathcal O_J(t_{\rm dim,\star})\big|_{\mathcal H_J}
  \longrightarrow
  \begin{pmatrix}
  0&\sqrt J\\
  \sqrt J&-J
  \end{pmatrix}.
  \label{eq:dimensional-devries-target}
\end{equation}
The positive branch is then assigned to the pole-ratio target, and the negative
branch is assigned to a scalar, boundary, brane, or compactification datum.
The compatibility diagram is
\begin{equation}
\begin{aligned}
(D,n_{\rm KK})=(11,7)\ {\rm or}\ (7,3)
&\xrightarrow{\ t_{\rm dim}\ }
(D,n_{\rm KK})=(9,5)\ {\rm or}\ (5,1),
\\
(D,n_{\rm KK})=(10,6)\ {\rm or}\ (6,2)
&:\quad
\mathcal O_J(t_{\rm dim,\star})\big|_{\mathcal H_J}\to K_J(\lambda),
\\
K_J(\lambda)
&\longrightarrow
(J_W,J_Z)=(J_H,J_{\rm adj}),
\\
(J_W,J_Z)
&\longrightarrow
\Delta^{-1}_{V,T}(s_{W,Z}) .
\end{aligned}
\label{eq:dimensional-compatibility-diagram}
\end{equation}
Thus the dimensional interpolation can contribute to the paper only through a
kernel derivation, a map from dimensional data to the electroweak ray, an
ordered assignment theorem, and a pole-matching theorem.

\subsection{Top-quark scale}
\label{app:top-quark-scale}

The dimensional interpolation also creates a top-sector question.  In the
Standard Model, fermion masses are generated by Yukawa couplings to the Higgs
vacuum expectation value.  The Higgs coupling to a fermion is proportional to
its mass, so the top quark has the largest direct fermionic coupling to the
electroweak order parameter.  This fact is visible in Higgs phenomenology and in
the role of top loops in Higgs production, electroweak corrections, and vacuum
stability \cite{Logan2022,Dawson2018,AlvarezGaumeVazquezMozo2023}.

The DeVries question is sharper: does the six-dimensional interior sector
select the top quark as the fermion attached most directly to the order
parameter?  A positive answer would require a source-theory map
\begin{equation}
  \mathcal F_{\rm top}:
  \mathcal H_J
  \longrightarrow
  y_t,\quad m_t,\quad
  \hbox{or a top-sector boundary datum},
  \label{eq:top-sector-map}
\end{equation}
compatible with the same determinant that gives the W/Z pole-ratio clue.  This
is an open target.  The present manuscript records the source-backed Standard
Model fact and the required form of a DeVries explanation.
